% ======================================================================
%  REVISÃO FINAL P1 — Erros óbvios que o professor usa como armadilha
%  Leia isso 30 minutos antes da prova. Nada mais.
% ======================================================================
\documentclass[10pt, a4paper]{article}
\usepackage[utf8]{inputenc}
\usepackage[T1]{fontenc}
\usepackage{lmodern}
\usepackage[brazil]{babel}
\usepackage[top=1.5cm, bottom=1.5cm, left=2cm, right=2cm]{geometry}
\usepackage{amsmath, amssymb}
\usepackage{xcolor}
\usepackage{tcolorbox}
\tcbuselibrary{breakable, skins}
\usepackage{enumitem}
\usepackage{multicol}
\usepackage{booktabs}
\usepackage{microtype}
\usepackage{setspace}

\definecolor{vermelho}{HTML}{991b1b}
\definecolor{verde}{HTML}{166534}
\definecolor{azul}{HTML}{1e3a5f}
\definecolor{laranja}{HTML}{92400e}
\definecolor{cinza}{HTML}{374151}

\newtcolorbox{errobox}[1]{standard, colback=white, colframe=vermelho,
  boxrule=1pt, arc=2pt, left=6pt, right=6pt, top=4pt, bottom=4pt,
  fonttitle=\bfseries\small\color{vermelho}, title={#1}}

\newtcolorbox{certobox}[1]{standard, colback=white, colframe=verde,
  boxrule=0.8pt, arc=2pt, left=6pt, right=6pt, top=4pt, bottom=4pt,
  fonttitle=\bfseries\small\color{verde}, title={#1}}

\newtcolorbox{formulabox}{standard, colback=white, colframe=azul,
  boxrule=0.8pt, arc=2pt, left=6pt, right=6pt, top=4pt, bottom=4pt}

\setlength{\parskip}{0pt}
\setlength{\parindent}{0pt}
\setlist[itemize]{nosep, leftmargin=*, label=\textbf{--}}
\setlist[enumerate]{nosep, leftmargin=*, label=\textbf{\arabic*.}}

\pagestyle{empty}

\begin{document}

\begin{center}
{\large\bfseries\color{azul} REVISÃO FINAL P1 — MACROECONOMIA I}\\
{\small\color{cinza} Apenas erros óbvios e armadilhas confirmadas da P1 2025}
\end{center}

\vspace{0.3em}
\begin{multicols}{2}

% =====================================================================
\textbf{\color{azul}\large 1. CRRA e $\theta$}

\begin{errobox}{FALSO — o professor usou isso na P1 2025 Q1}
``$\theta$ é a elasticidade de substituição intertemporal.''
\end{errobox}
\vspace{0.3em}
\begin{certobox}{CERTO}
$\theta$ é o inverso da EIS:\quad $\displaystyle\text{EIS} = \frac{1}{\theta}$\\
$\theta\uparrow$ $\Rightarrow$ consumidor \emph{menos} responsivo a $r$.
\end{certobox}

\begin{errobox}{FALSO — P1 2025 Q1}
``$u' = \theta C^{-\theta}$''
\end{errobox}
\vspace{0.3em}
\begin{certobox}{CERTO}
$u' = C^{-\theta}$ \quad(sem o $\theta$ multiplicando)
\end{certobox}

\vspace{0.5em}
% =====================================================================
\textbf{\color{azul}\large 2. Equação de Euler}

\begin{formulabox}
\[\frac{\dot C}{C} = \frac{r - \rho}{\theta}\]
\end{formulabox}

\begin{errobox}{FALSO — P1 2025 Q2}
``A taxa de desconto $\rho$ afeta \emph{positivamente} a taxa de crescimento do consumo.''
\end{errobox}
\vspace{0.3em}
\begin{certobox}{CERTO}
$\rho$ tem sinal negativo: $\rho\uparrow$ \emph{reduz} $\dot C/C$.\\
Mais impaciência → consumir hoje, crescer menos.
\end{certobox}

\begin{errobox}{FALSO}
``A equação de Euler implica que famílias \emph{diminuem} o crescimento do consumo frente a elevação de juros.''
\end{errobox}
\vspace{0.3em}
\begin{certobox}{CERTO}
$r\uparrow$ $\Rightarrow$ $\dot C/C\uparrow$ — juros maiores estimulam poupança e crescimento do consumo.
\end{certobox}

\vspace{0.5em}
% =====================================================================
\textbf{\color{azul}\large 3. Diagrama de Fase RCK}

\begin{certobox}{Quadrantes — decorar}
\begin{tabular}{lcc}
\toprule
Posição & $\dot c$ & $\dot k$ \\
\midrule
Esq.$\dot c$, Acima $\dot k$ & $+$ & $-$ \\
Esq.$\dot c$, Abaixo $\dot k$ & $+$ & $+$ \\
Dir.$\dot c$, Acima $\dot k$ & $-$ & $-$ \\
\textbf{Dir.$\dot c$, Abaixo $\dot k$} & $-$ & $+$ \\
\bottomrule
\end{tabular}\\
\small A última linha = gabarito Q4 P1-2025.
\end{certobox}

\begin{errobox}{FALSO}
``O crescimento equilibrado do RCK ocorre na Regra de Ouro.''
\end{errobox}
\vspace{0.3em}
\begin{certobox}{CERTO}
RCK: $f'(k^*) = \delta+\rho+\theta g$ \quad(não é Regra de Ouro)\\
RO: $f'(k^{RO}) = \delta+n+g$ \quad(maximiza consumo no SS)\\
Em geral $k^* < k^{RO}$.
\end{certobox}

\begin{errobox}{FALSO}
``O crescimento populacional está \emph{positivamente} relacionado com a variação do capital.''
\end{errobox}
\vspace{0.3em}
\begin{certobox}{CERTO}
$\dot k = f(k)-c-(n+g+\delta)k$: $n$ aparece com sinal negativo.
\end{certobox}

\columnbreak
% =====================================================================
\textbf{\color{azul}\large 4. Choque Tecnológico RBC}

\begin{certobox}{Os 3 canais (Lista I Q8, Q7 P1-2025)}
\textbf{A.} Efeito riqueza: $C\uparrow$, $L\downarrow$\\
\textbf{B.} Sub.\ intertemporal: $I\uparrow$, $L$ corrente $\uparrow$\\
\textbf{C.} Produtividade: $w\uparrow$ $\Rightarrow$ $L\uparrow$
\end{certobox}

\begin{errobox}{FALSO — P1 2025 Q7}
``Choque positivo estimula \emph{redução} da poupança.''\\
``Choque positivo: salários sobem mas taxa de juros \emph{não} sobe.''\\
``Maior persistência → oferta de $L$ fica \emph{abaixo} do nível inicial.''
\end{errobox}
\vspace{0.3em}
\begin{certobox}{CERTO}
$I\uparrow$ = poupança sobe.\\
$f'(K)\uparrow$ = taxa de retorno do capital também sobe.\\
Maior persistência → efeito sub.\ intertemporal mais forte → $L$ corrente \emph{acima} do nível inicial.
\end{certobox}

\vspace{0.5em}
% =====================================================================
\textbf{\color{azul}\large 5. Mark-up e Salário Real}

\begin{formulabox}
$P = \mu \cdot W/f'(L)$ \quad$\Rightarrow$\quad $W/P = f'(L)/\mu$
\end{formulabox}

\begin{certobox}{Gabarito Q8 P1-2025}
Mark-up contra-cíclico ($\mu\downarrow$ quando $Y\uparrow$) + $W$ fixo\\
$\Rightarrow$ $W/P\uparrow$ quando $Y\uparrow$ = salário real \textbf{pró-cíclico}.
\end{certobox}

\begin{errobox}{FALSO}
``Se $W$ e $\mu$ são constantes, choque negativo de demanda \emph{reduz} salário real.''
\end{errobox}
\vspace{0.3em}
\begin{certobox}{CERTO}
Se ambos constantes → $W/P$ não muda. Ponto.
\end{certobox}

\vspace{0.5em}
% =====================================================================
\textbf{\color{azul}\large 6. Lucas: Sinal e Crítica}

\begin{errobox}{FALSO — P1 2025 Q11}
``Crítica de Lucas: previsões devem assumir \emph{manutenção} do estado de expectativas diante de mudança de regime.''
\end{errobox}
\vspace{0.3em}
\begin{certobox}{CERTO}
A crítica diz o oposto: quando a regra muda, os agentes \emph{revisam} expectativas e os parâmetros estimados se alteram. Modelo antigo fica inválido.
\end{certobox}

\begin{errobox}{FALSO — P1 2025 Q11}
``Ausência de trade-off LP está associada a $b > 1$.''
\end{errobox}
\vspace{0.3em}
\begin{certobox}{CERTO}
Decorre de $p = E[p]$ no LP — independe do valor de $b$.
\end{certobox}

\vspace{0.5em}
% =====================================================================
\textbf{\color{azul}\large 7. Fischer vs.\ Taylor}

\begin{errobox}{FALSO — P1 2025 Q13}
``A diferença básica está na formação de preços: Fischer usa contratos, Taylor não.''
\end{errobox}
\vspace{0.3em}
\begin{certobox}{CERTO}
\emph{Ambos} usam contratos escalonados. Diferença = timing:\\
Fischer: expectativas \textbf{defasadas} ($E_{t-2}$, $E_{t-1}$)\\
Taylor: expectativas \textbf{forward-looking} racionais
\end{certobox}

\begin{errobox}{FALSO — P1 2025 Q13}
``Fischer com $\phi$ elevado: efeitos monetários sobre $Y$ \emph{maiores}.''
\end{errobox}
\vspace{0.3em}
\begin{certobox}{CERTO}
Multiplicador $= 1/(1+\phi)$: $\phi\uparrow$ $\Rightarrow$ multiplicador \emph{menor}.\\
Taylor: $\lambda = (1-\phi)/(1+\phi)$: $\phi\uparrow$ $\Rightarrow$ $\lambda\downarrow$ (menos persistência).
\end{certobox}

\vspace{0.5em}
% =====================================================================
\textbf{\color{azul}\large 8. Calvo: $\alpha$, $\kappa$, $\phi$}

\begin{formulabox}
$\kappa = \dfrac{\alpha[1-(1-\alpha)\beta]\phi}{1-\alpha}$ \qquad $\phi = \theta+\gamma-1$
\end{formulabox}

\begin{errobox}{ATENÇÃO — inversão clássica}
$\alpha\uparrow$ (mais firmas presas) $\Rightarrow$ $\kappa\downarrow$ (curva mais plana)\\
Plana = inflação surda ao hiato = sacrifício alto na desinflação.
\end{errobox}
\vspace{0.3em}
\begin{certobox}{CERTO — $\phi\uparrow$}
$\phi\uparrow$ $\Rightarrow$ custo marginal mais sensível ao hiato $\Rightarrow$ $\kappa\uparrow$\\
Firmas que reajustam pedem preços muito maiores $\Rightarrow$ inflação reage mais.
\end{certobox}

\begin{errobox}{FALSO — P1 2025 Q14}
``Calvo pressupõe firmas \emph{heterogêneas}.''\\
``Calvo tem componente backward-looking.''
\end{errobox}
\vspace{0.3em}
\begin{certobox}{CERTO}
Firmas idênticas ex ante; heterogeneidade só em quem conseguiu reajustar.\\
Backward-looking = CEE (2005), \emph{não} Calvo puro.
\end{certobox}

\vspace{0.5em}
% =====================================================================
\textbf{\color{azul}\large 9. $p^* = \phi m + (1-\phi)p$ e Phillips Híbrida}

\begin{errobox}{FALSO — P1 2025 Q12}
``Maior $\phi$ $\Rightarrow$ maior rigidez.''
\end{errobox}
\vspace{0.3em}
\begin{certobox}{CERTO — gabarito = d}
$\phi\uparrow$ $\Rightarrow$ \emph{menos} rigidez. $\phi=1$ $\Rightarrow$ preços totalmente flexíveis = RBC.
\end{certobox}

\begin{errobox}{FALSO — P1 2025 Q9}
``$\varepsilon_t^s$ é relacionado à gestão monetária do BC.''
\end{errobox}
\vspace{0.3em}
\begin{certobox}{CERTO}
$\varepsilon_t^s$ = choque de \emph{oferta} (supply shock). Nada a ver com política monetária.
\end{certobox}

\begin{certobox}{Phillips Híbrida — condição-chave}
$\pi_t = \gamma\pi_{t-1} + (1-\gamma)E_t\pi_{t+1} + \chi y_t$\\
$\gamma > 1/2$ $\Rightarrow$ inércia domina $\Rightarrow$ razão de sacrifício \textbf{positiva}\\
$\gamma < 1/2$ $\Rightarrow$ forward-looking domina $\Rightarrow$ desinflação mais barata
\end{certobox}

\vspace{0.5em}
% =====================================================================
\textbf{\color{azul}\large 10. Oferta de Trabalho: $U = C - \frac{1}{d}L^d$}

\begin{formulabox}
CPO: $W/P - L^{d-1} = 0$ $\Rightarrow$ $L^* = (W/P)^{1/(d-1)}$
\end{formulabox}

\begin{errobox}{FALSO — P1 2025 Q10}
``Relação \emph{inversa} entre salário real e oferta de trabalho.''\\
``Utilidade marginal do trabalho é \emph{positiva}.''
\end{errobox}
\vspace{0.3em}
\begin{certobox}{CERTO}
Relação \emph{positiva}: $W/P\uparrow$ $\Rightarrow$ $L^*\uparrow$.\\
No ótimo, $\partial U/\partial L = 0$ (se fosse positiva, valeria trabalhar mais).
\end{certobox}

\end{multicols}

\vspace{0.5em}
\begin{tcolorbox}[standard, colback=white, colframe=azul, boxrule=1.2pt,
  arc=3pt, left=8pt, right=8pt, top=5pt, bottom=5pt]
\textbf{\color{azul} Regra de ouro antes de marcar:} Se a alternativa tem uma relação positiva onde você esperaria negativa (ou vice-versa), ou inverte o papel de um parâmetro, provavelmente é a armadilha. O professor usa \emph{sempre} o mesmo padrão: inversão de sinal, confusão de parâmetros ($\theta$ vs EIS; $\alpha$ vs $\kappa$; $\phi$ vs rigidez) e afirmativas com ``aumenta'' onde o correto é ``diminui''.
\end{tcolorbox}

\end{document}
